\section{Discussion and Limitations}
The expanded design links construct identity, nutrient-sensitive distribution, selected regulatory sites and downstream function. Its strength comes from comparisons that separate competing explanations: matched nutrient/solvent controls, endogenous confirmation, expression titration, individual and combined site substitutions, genetic perturbation, and loss-and-rescue. A planned sequence of execution is construct acceptance, cell-assay pilot, localization/recovery confirmation, then mechanism and functional modules. Failed early quality checks are resolved before downstream interpretation.

Each control has a defined scope. EGFP-only does not prove that a TFEB fusion behaves exactly like endogenous protein. Site substitutions perturb local chemistry as well as phosphorylation potential. RPTOR/RICTOR comparisons strengthen complex attribution but require competent knockdown and may reveal indirect effects; direct phosphorylation requires a separate biochemical experiment. Recovery of N/C is net redistribution, not an independently measured export rate. TFEB knockdown and rescue can support a contribution without establishing that TFEB is the only MiT/TFE-family mediator.

Nuclear accumulation, transcription and autophagic turnover are distinct endpoints. A stress response can move TFEB into nuclei without improving degradation. Drug-induced autophagy can also involve TFEB-independent pathways. Consequently, the proposed functional conclusions require their own genetic and inhibitor controls. If those assays disagree, the conclusion is narrowed to the supported level rather than describing every endpoint as ``activation.''

The C-terminal extension adds a question beyond reproducing redistribution. Its factorial design separates tag-dependent and stimulus-dependent effects, while 3T supplies a complementary chemical perturbation. Neither alanine nor threonine substitution isolates phosphorylation from all structural consequences. Matched nuclear exposure, target engagement and phosphosite specificity are therefore necessary to interpret transcriptional differences. A null result with a wide interval remains inconclusive. The study does not claim that this exact experimental combination is unprecedented.

This remains an experimental proposal: no wet-lab clones, images, blots, qPCR measurements or flux data have been acquired. The new sequence files establish intended mutant identities, not their cellular activity. Supplier-specific reagents, siRNA targets, qPCR oligos, antibody performance and final doses/sample size require documented pilot validation. These implementation items and actual author contributions must be resolved before execution or course submission. The design can address major foreseeable objections; it cannot guarantee that every alternative mechanism has been excluded.
